% VecDB@VLDB 2026 workshop format -- uses the shipped acmart.cls + ACM-Reference-Format.bst
% Tables are generated from JSON artefacts by
%   `python -m ebrag.benchmarks.experiments.c4_render_tables`

\documentclass[sigconf, nonacm, natbib]{acmart}

\let\Bbbk\relax       % avoid amssymb / amsfonts conflict pulled in by acmart
\usepackage{amssymb}  % \checkmark in the per-model ladder table

%%% VLDB workshop metadata (required) %%%
\newcommand\vldbyear{2026}
\newcommand\vldbworkshop{The 2nd Workshop on Vector Databases}
\newcommand\vldbauthors{\authors}
\newcommand\vldbtitle{\shorttitle}
\newcommand\vldbavailabilityurl{https://github.com/sarkar-dipankar/ebrag-vecdb-2026-paper}
\newcommand\vldbpagestyle{plain}

\begin{document}

% Relax line-breaking penalties to suppress minor overfull-hbox warnings on
% hyphenated compounds in this two-column format.
\emergencystretch=3em
\hyphenpenalty=100
\tolerance=1000

\title{An Input-Regime Audit of Conflict Detection for Retrieval-Augmented Generation}

\author{Dipankar Sarkar}
\affiliation{%
  \institution{Skelf Research}
  \country{}
}
\email{dipankar@skelfresearch.com}

\begin{abstract}
Vector-database-backed retrieval-augmented generation (RAG) pipelines increasingly add a
\emph{conflict-detection} stage that flags contradictions among retrieved passages before
generation. Two recent directions dominate practice. The first is a pairwise NLI
cross-encoder over retrieved passages, the standard baseline detector across
conflict-aware RAG systems and the stack deployed by
\textsc{ConflictRAG}~\citep{conflictrag2025}. The second is a
\emph{query-conditioned} formulation that lets the NLI label depend on the question
being answered, recently benchmarked by~\citet{canby2025qnli}. Practitioners must choose
between them with little empirical guidance. We provide that guidance. We audit three
detectors: the classic DeBERTa-MNLI cross-encoder, an LLM-as-judge baseline, and the
same LLM judge question-conditioned. We evaluate them on curated probes ($n{=}63$ with
bootstrap CIs), a 10-model LLM sweep, a prompt-robustness ablation, a long-passage
realism step, an answer-bearing scoping study, and external validation on 100
human-labelled FEVER pairs. Five findings guide deployment. (1) LLM judges retain
perfect precision on labelled FEVER pairs and flag far fewer top-$k$ retrieved non-gold
passages than the cross-encoder ($0.09$ vs $0.43$ flag rate).
(2) A same-name ``sense-trap'' false positive
systematically fools both the cross-encoder and pairwise LLM judging. Only
question-conditioning closes the gap (precision $0.68 \to 1.00$ on curated probes,
replicated for 8 of 10 LLM families). (3) The cross-encoder vs LLM trade is
\emph{input-distribution dependent}. The cross-encoder is competitive on its native
short claim/evidence distribution (FEVER F1 $0.85$) but unusable on long multi-claim
passages (F1 $0.18$). (4) The small-model regression under the question-conditioning
prompt is a capability limit, not a format limit. We verify this by inspecting raw
outputs. (5) Pairwise conflict detection is ill-posed on multi-hop evidence fragments
and must be scoped to answer-bearing units. We synthesise the findings into a five-row
input-regime ladder. The ladder recommends a detector for each retrieval context a
vector-DB practitioner is likely to encounter. Code and result artefacts are released.
\end{abstract}

\maketitle

%%% do not modify the following VLDB block %%
%%% VLDB block start %%%
\pagestyle{\vldbpagestyle}
\begingroup\small\noindent\raggedright\textbf{VLDB Workshop Reference Format:}\\
\vldbauthors. \vldbtitle. VLDB \vldbyear\ Workshop: \vldbworkshop.\\
\endgroup
\begingroup
\renewcommand\thefootnote{}\footnote{\noindent
This work is licensed under the Creative Commons BY-NC-ND 4.0 International License. Visit \url{https://creativecommons.org/licenses/by-nc-nd/4.0/} to view a copy of this license. For any use beyond those covered by this license, obtain permission by emailing \href{mailto:info@vldb.org}{info@vldb.org}. Copyright is held by the owner/author(s). Publication rights licensed to the VLDB Endowment. \\
\raggedright Proceedings of the VLDB Endowment. ISSN 2150-8097. \\
}\addtocounter{footnote}{-1}\endgroup
%%% VLDB block end %%%

%%% do not modify the following VLDB block %%
%%% VLDB block start %%%
\ifdefempty{\vldbavailabilityurl}{}{
\vspace{.3cm}
\begingroup\small\noindent\raggedright\textbf{VLDB Workshop Artifact Availability:}\\
The source code, data, and/or other artifacts have been made available at \url{\vldbavailabilityurl}.
\endgroup
}
%%% VLDB block end %%%

\section{Introduction}
\label{sec:intro}
A vector-database-backed RAG pipeline returns a top-$k$ set of passages. A growing class
of \emph{conflict-aware} systems then runs a \emph{conflict-detection} stage on those
passages before generation~\citep{madamrag2025,astuterag2024,contradetect2025,
fvarag2025,clasheval2024,worsezero2025}. The detector's signal gates downstream choices:
multi-agent debate over conflicting sources~\citep{madamrag2025}, source-reliability
consolidation~\citep{astuterag2024}, structured support/refute
output~\citep{conrag2025}, or calibrated abstention via standard confidence
calibration~\citep{guo2017calibration}. The detector is the load-bearing component
upstream of all of these: a poor detector silently caps every method built on top of it.

Two recent directions dominate practice. A pairwise NLI cross-encoder applied to
retrieved-passage pairs is the standard baseline detector across conflict-aware
pipelines~\citep{conflictrag2025,madamrag2025,contradetect2025}. As a recent
conflict-aware RAG system, \textsc{ConflictRAG}~\citep{conflictrag2025} reports
$88.7\%$ F1 and downstream correctness gains of $5.3$ to $6.1\%$ over conflict-aware
baselines.
Concurrent work by~\citet{canby2025qnli} formalises and benchmarks \emph{query-conditioned
NLI}: the entailment label depends on a premise, a hypothesis, \emph{and} a query that
specifies the aspect under comparison. A practitioner picking a detector for a
vector-DB pipeline needs to know three things: which approach, on what kind of input,
and with what model.

We provide that empirical guidance. We audit a representative pairwise NLI cross-encoder
(DeBERTa-base-MNLI), an LLM-as-judge baseline in the style
of~\citet{contradetect2025}, and the same LLM judge question-conditioned in the style
of~\citet{canby2025qnli}. The setting is the conflict-aware RAG pipeline of
\textsc{ConflictRAG}~\citep{conflictrag2025} and adjacent systems. We evaluate the three
detectors across a curated probe set, a 10-model LLM sweep, a small-model
prompt-robustness ablation, a long-passage realism study, an answer-bearing scoping
experiment, and external validation on 100 human-labelled FEVER claim/evidence pairs.
We report bootstrap 95\% confidence intervals throughout, so the audit's orderings are
statistically interpretable.

\paragraph{Where this sits in the vector-DB pipeline}
The detector we audit is a single stage in a longer vector-DB-backed RAG pipeline:
\begin{center}\small
top-$k$ retrieval $\to$ passage units $\to$ pair construction $\to$\\
detector $\to$ generation policy
\end{center}
The vector-DB practitioner controls five design choices that change what the detector
sees: (i)~chunk size and granularity (sentence, paragraph, multi-hop); (ii)~top-$k$ and
how many pairs are checked; (iii)~whether pairs are formed sentence-level,
paragraph-level, or only on \emph{answer-bearing} units; (iv)~whether conflict detection
runs before or after a question-conditioned extraction step; and (v)~detector cost and
latency vs precision. Our audit holds the surrounding pipeline fixed and varies the
detector. The findings below land on choices (i), (iii), (iv), and (v).

\paragraph{Contributions}
\begin{enumerate}
  \item \textbf{A mechanistic failure mode of pairwise judging: the sense trap.}
  Same-name-different-referent passage pairs (``Paris is the capital of France'' vs.\
  ``Paris is a town in Texas'') systematically fool both the pairwise NLI cross-encoder
  baseline used across conflict-aware RAG
  \citep{conflictrag2025,madamrag2025,contradetect2025} and a pairwise LLM judge.
  We isolate the mechanism: pairwise judging lacks the question that says
  \emph{which} referent is relevant.
  \item \textbf{An empirical audit of the query-conditioning fix.}
  Conditioning the LLM judge on the question is the operational form of the
  query-conditioned NLI proposed by~\citet{canby2025qnli}. It closes the sense-trap gap
  on the curated set (precision $0.68 \to 1.00$) and replicates across 8 of 10
  open-weight LLM families (\S\ref{sec:sweep}) with non-overlapping bootstrap CIs.
  \item \textbf{Capability, not format.} The small-Gemma regression under the
  question-conditioning prompt is a reasoning limit, not a JSON-format artefact.
  Raw model output inspection shows it (\S\ref{sec:promptrob}).
  \item \textbf{Scoping: pairwise conflict is ill-posed on multi-hop evidence
  fragments.} An evidence paragraph about one entity (``Scott Derrickson is American'')
  does not contradict a relationship verdict (``Derrickson and Wood are of different
  nationalities''). The judge's \texttt{neutral} is correct. The audit must be scoped
  to answer-bearing pairs. This is a methodological clarification useful to any
  downstream vector-DB pipeline that retrieves multi-hop evidence.
  \item \textbf{The detector trade is input-regime dependent.} The cross-encoder is
  competitive on its native short claim/evidence distribution (FEVER F1 $0.85$) but
  unusable on long multi-claim passages (F1 $0.18$). LLM judges are robust across
  lengths but more conservative on recall. We chart the five-row regime
  ladder (Table~\ref{tab:synthesis}) as a deployment guide.
\end{enumerate}

\paragraph{Relation to prior work}
\citet{canby2025qnli} formalise query-conditioned NLI as a benchmark task; we instead
audit it as a \emph{deployable conflict-detection signal} for RAG. We characterise where
its fix dominates the pairwise alternative and identify the model-capability regime in
which it backfires. \citet{conflictrag2025} build a two-stage conflict-aware RAG pipeline
(embedding-based MLP classifier with selective LLM refinement, then source-credibility
weighting) and report end-task gains; we audit the underlying pairwise NLI detector
signal in isolation, upstream of any source-credibility stage. \citet{conqret2025} and
\citet{madamrag2025} target the generator's handling of conflicting evidence; we work on
the detector that supplies their input signal.

\section{Related Work}
\label{sec:related}

\paragraph{Conflict detection on top of retrieval}
\textsc{ConflictRAG}~\citep{conflictrag2025} is the most direct system motivation for
our audit. It introduces a two-stage conflict-detection module (an embedding-based MLP
classifier with selective LLM refinement) followed by source-credibility weighting,
reporting $88.7\%$ conflict-detection F1 and $5.3$ to $6.1\%$ downstream correctness
gains over conflict-aware baselines. Our audit complements that system result by
isolating the pairwise NLI signal, the standard cross-encoder baseline used across this
literature~\citep{madamrag2025,contradetect2025}, and characterising where it succeeds
and where it fails. \citet{contradetect2025} evaluate LLMs as context-validators in RAG
pipelines, a direct counterpart to our \texttt{llm\_bare} detector.
\citet{worsezero2025} construct a fact-checking dataset for RAG robustness against
misleading retrieval, complementing the FEVER claim/evidence audit we run in
\S\ref{sec:fever}.

\paragraph{Query-conditioned NLI}
\citet{canby2025qnli} introduce and benchmark query-conditioned NLI. The
entailment/contradiction label is conditioned on a query alongside the premise and
hypothesis. Their contribution is the formalisation and a benchmark. Ours is an
\emph{empirical deployment audit} of the same idea for RAG. We ask how the
question-conditioned variant behaves vs.\ pairwise NLI across model families
(\S\ref{sec:sweep}), input lengths (\S\ref{sec:realism}, \S\ref{sec:answerbearing}), and
real human-labelled retrieved evidence (\S\ref{sec:fever}). We also identify the
small-model capability regime in which the fix regresses (\S\ref{sec:promptrob}).

\paragraph{Generator-side conflict-aware RAG}
A line of work targets how the \emph{generator} handles disagreement: multi-agent
grounded debate (\textsc{Madam-Rag} + \textsc{RAMDocs})~\citep{madamrag2025};
source-reliability consolidation~\citep{astuterag2024}; structured support/refute output
\citep{conrag2025}; falsification for sycophancy~\citep{fvarag2025}; the
prior-vs-context tug-of-war \citep{clasheval2024}; stance-diversified retrieved evidence
for argumentation \citep{conqret2025}. All of these consume a detector signal and
reason over conflicting inputs; none audit the detector itself. We operate strictly
upstream, measuring where pairwise detection (cross-encoder NLI and LLM-as-judge alike)
fails and where question-conditioning closes the gap, so the downstream choices these
systems make rest on a characterised signal.

\paragraph{NLI as a building block}
Cross-encoder NLI predates the RAG conflict-detection literature. We use
DeBERTa-base-MNLI as the canonical pairwise NLI representative because variants of this
model family are the standard baseline detector across conflict-aware
RAG~\citep{conflictrag2025,madamrag2025,contradetect2025}. Temperature scaling and
calibration of NLI heads \citep{guo2017calibration} are orthogonal to the audit. We
report contradiction-class probabilities directly.

\section{Detectors}
\label{sec:detectors}
We compare three detectors. All are evaluated as a binary contradiction classifier.
\begin{description}
  \item[\textsc{cross\_encoder\_nli}.] DeBERTa-base-MNLI, the ``classic'' pairwise
  detector used by many conflict-aware RAG pipelines. Input: a single text-pair string.
  Pairwise by design, with no slot for the question.
  \item[\textsc{llm\_bare}.] An LLM judge prompted to label the pair entailment /
  contradiction / neutral. Pairwise like the cross-encoder, but with arbitrary
  instruction-following.
  \item[\textsc{llm\_answer\_relevant}.] The same LLM judge prompted to decide whether
  passages A and B give \emph{conflicting answers to a specific question}. The question
  enters the prompt; A and B are otherwise the same texts as in the bare variant. This
  is the operational form of the query-conditioned NLI of~\citet{canby2025qnli}, used
  here as a deployable conflict-detection signal.
\end{description}

\section{Probes and data}
\label{sec:data}
\textbf{Curated probe set.} 63 probes spanning 9 categories: same-name sense traps,
explicit conflicts, numeric/temporal/implicit conflicts, paraphrase agreement, unit
equivalence, negation agreement, topic shift. Gold-correctness invariants are unit-tested
(every conflict category is gold=conflict; every no-conflict category is
gold=no-conflict; template-generated numeric/temporal/unit pairs differ in only the
controlled value).

\textbf{Long answer-bearing pairs.} For 14 single-hop factoid questions with two
competing answers, we LLM-write 3 to 4 sentence passages that each assert a specific
answer. This gives conflict, agreement, and topic-shift pairs over long real-sounding
text.

\textbf{FEVER (real, human-labelled).} 100 VERIFIABLE pairs from
\texttt{copenlu/fever\_gold\_evidence}: 50 REFUTES (conflict) and 50 SUPPORTS
(no-conflict). Premise is the gold evidence sentence and hypothesis is the claim.

\section{The sense-trap and answer-relevance}
\label{sec:answerrel}
On the curated probes, a same-name-different-referent pair (``Paris, France'' /
``Paris, Texas'' with question \emph{What is the capital of France?}) is consistently
flagged as a contradiction by the cross-encoder and frequently by the bare LLM judge.
Both decide pairwise without the question and default to the assumption that the two
passages are about the same ``Paris''. Conditioning the judge on the question
(``Do A and B give conflicting answers to \emph{this} question?'') removes the failure
without recall loss on the genuine conflicts: precision rises from $0.68$ to $1.00$ on
the curated probes (Table~\ref{tab:detector_ladder}), and the cross-encoder's upper CI
($0.83$) sits below answer-relevance's point estimate ($1.00$).

\begin{table}[t]
\centering
\small
\begin{tabular}{lccc}
\toprule
Detector & Precision (95\% CI) & F1 & Sense-trap acc (95\% CI) \\
\midrule
cross-enc NLI & 0.68 [0.52,0.83] & 0.81 & 0.67 [0.38,0.92] \\
LLM bare & 0.90 [0.77,1.00] & 0.95 & 0.75 [0.46,1.00] \\
LLM answer-rel & 1.00 [1.00,1.00] & 1.00 & 1.00 [1.00,1.00] \\
\bottomrule
\end{tabular}
\caption{Detector ladder on the curated probe set ($n=63$). The cross-encoder is structurally
unable to be question-conditioned; only the answer-relevant LLM judge closes the sense-trap gap.}
\label{tab:detector_ladder}
\end{table}


\section{Cross-model scaling}
\label{sec:sweep}
We run the bare vs.~answer-relevant comparison across the 10-model
\texttt{DEFAULT\_SWEEP} (six families, 4B to 671B parameters) with bootstrap precision
CIs (Table~\ref{tab:sweep}). Across families, the four most capable models
(gpt-oss:20b, gpt-oss:120b, qwen3-next:80b, glm-4.6) reach perfect precision under the
answer-relevance prompt with bare CIs that sit clearly below. 8 of 10 models improve.
The two small Gemma models regress with non-overlapping CIs: gemma3:4b collapses to
P $0.51$ [$0.37, 0.65$] (FP $1{\to}25$). The mean lift is small ($+0.011$) because the
two regressions drag it down; the headline is the family-spanning capability
dependence, not the mean.

\begin{table*}[t]
\centering
\small
\begin{tabular}{llrcccc}
\toprule
Model & Family & Size (B) & bare P (95\% CI) & answer-rel P (95\% CI) & Lift & FP b$\rightarrow$ar \\
\midrule
\texttt{gpt-oss:20b} & gpt-oss & 20 & 0.84 [0.70,0.96] & 1.00 [1.00,1.00] & +0.16 & 5$\rightarrow$0 \\
\texttt{gpt-oss:120b} & gpt-oss & 120 & 0.90 [0.77,1.00] & 1.00 [1.00,1.00] & +0.10 & 3$\rightarrow$0 \\
\texttt{gemma3:4b} & gemma & 4 & 0.96 [0.88,1.00] & 0.51 [0.37,0.65] & -0.45 & 1$\rightarrow$25 \\
\texttt{gemma3:12b} & gemma & 12 & 0.76 [0.62,0.90] & 0.63 [0.47,0.78] & -0.13 & 8$\rightarrow$15 \\
\texttt{gemma3:27b} & gemma & 27 & 0.83 [0.69,0.96] & 0.87 [0.73,0.97] & +0.03 & 5$\rightarrow$4 \\
\texttt{qwen3-next:80b} & qwen & 80 & 0.93 [0.81,1.00] & 1.00 [1.00,1.00] & +0.07 & 2$\rightarrow$0 \\
\texttt{deepseek-v3.1:671b} & deepseek & 671 & 0.84 [0.70,0.95] & 0.87 [0.73,0.97] & +0.03 & 5$\rightarrow$4 \\
\texttt{glm-4.6} & glm & -- & 0.84 [0.70,0.96] & 1.00 [1.00,1.00] & +0.16 & 5$\rightarrow$0 \\
\texttt{ministral-3:8b} & mistral & 8 & 0.70 [0.56,0.84] & 0.74 [0.59,0.88] & +0.04 & 11$\rightarrow$9 \\
\texttt{nemotron-3-nano:30b} & nemotron & 30 & 0.87 [0.73,0.97] & 0.96 [0.88,1.00] & +0.10 & 4$\rightarrow$1 \\
\bottomrule
\end{tabular}
\caption{Answer-relevance precision lift across the 10-model sweep
(8/10 improve; mean lift +0.011).
The four most capable models reach perfect precision; the two small Gemma models regress
(non-overlapping CIs). $n=63$ probes.}
\label{tab:sweep}
\end{table*}


\paragraph{The ladder against a fixed cross-encoder}
The cross-encoder NLI is model-fixed (DeBERTa-MNLI), so its precision on the curated set
is a constant reference. The per-LLM-model picture against it is in
Table~\ref{tab:ladder_per_model}. \textbf{Bare LLM judging beats the cross-encoder for
every one of the 10 LLM models} on this probe set. Answer-relevance beats it for 8 of
10, and the full ladder (answer-rel $\geq$ bare $\geq$ cross-encoder) holds for the
same 8 models that improved under answer-relevance in Table~\ref{tab:sweep}. The two
small Gemma regressors are the \emph{only} models on which an LLM judge would lose to
the DeBERTa-MNLI cross-encoder under the answer-relevant prompt: the binding constraint
is the capability limit (\S\ref{sec:promptrob}), not the framing.

\begin{table*}[t]
\centering
\small
\begin{tabular}{lcccc}
\toprule
LLM model & cross-enc P & LLM-bare P & LLM-ar P & ladder \\
\midrule
\texttt{gpt-oss:20b} & 0.68 & 0.84 & 1.00 & \checkmark \\
\texttt{gpt-oss:120b} & 0.68 & 0.90 & 1.00 & \checkmark \\
\texttt{gemma3:4b} & 0.68 & 0.96 & 0.51 & $\times$ \\
\texttt{gemma3:12b} & 0.68 & 0.76 & 0.63 & $\times$ \\
\texttt{gemma3:27b} & 0.68 & 0.83 & 0.87 & \checkmark \\
\texttt{qwen3-next:80b} & 0.68 & 0.93 & 1.00 & \checkmark \\
\texttt{deepseek-v3.1:671b} & 0.68 & 0.84 & 0.87 & \checkmark \\
\texttt{glm-4.6} & 0.68 & 0.84 & 1.00 & \checkmark \\
\texttt{ministral-3:8b} & 0.68 & 0.70 & 0.74 & \checkmark \\
\texttt{nemotron-3-nano:30b} & 0.68 & 0.87 & 0.96 & \checkmark \\
\bottomrule
\end{tabular}
\caption{Per-model ladder against the fixed cross-encoder reference
(DeBERTa-MNLI, $P=0.68$ on $n=63$). Bare LLM judging beats the cross-encoder for
\textbf{10/10} models; answer-relevance beats it for
\textbf{8/10}; the full ladder
(answer-rel $\geq$ bare $\geq$ cross-enc) holds for 8/10.
The two small Gemma regressors drop \emph{below} the cross-encoder under answer-relevance,
making them the only models where a cross-encoder would beat the LLM judge.}
\label{tab:ladder_per_model}
\end{table*}


\section{Capability, not format}
\label{sec:promptrob}
The small-Gemma regression could be a JSON-format failure or a reasoning limit. We
disentangle (Table~\ref{tab:prompt_robustness}) by judging under three answer-relevance
prompts: the production JSON prompt, a one-word \texttt{plain} variant, and a
\texttt{fewshot} variant with two worked examples. We test all three on gemma3:4b,
gemma3:12b, and the capable gpt-oss:20b control. We then \emph{inspect the raw outputs}
on the sense-trap false positives. gemma3:4b emits clean, parseable JSON that often
verbalises the right distinction. On Mercury planet vs.~element, its rationale reads
``Passage A answers the question, while Passage B discusses a chemical element''. The
model still labels the pair \texttt{contradiction}. This is a decision/reasoning
failure, not a format artefact. Switching to one-word output only shifts the threshold
(FP $25 \to 11$), and the model still fails on Paris/Cambridge sense traps. Few-shot
demonstration partially rescues gemma3:12b (sense-trap acc $0.50 \to 0.67$) but not
gemma3:4b. gpt-oss:20b is robust across all variants and goes to perfect with few-shot.

\begin{table*}[t]
\centering
\small
\begin{tabular}{llcccc}
\toprule
Model & Variant & Precision & F1 & Sense-trap acc & FP \\
\midrule
\texttt{gemma3:4b} & json & 0.51 & 0.68 & 0.00 & 25 \\
\texttt{gemma3:4b} & plain & 0.67 & 0.75 & 0.58 & 11 \\
\texttt{gemma3:4b} & fewshot & 0.62 & 0.76 & 0.42 & 16 \\
\texttt{gemma3:12b} & json & 0.63 & 0.78 & 0.50 & 15 \\
\texttt{gemma3:12b} & plain & 0.63 & 0.78 & 0.33 & 15 \\
\texttt{gemma3:12b} & fewshot & 0.76 & 0.87 & 0.67 & 8 \\
\texttt{gpt-oss:20b} & json & 0.96 & 0.98 & 0.92 & 1 \\
\texttt{gpt-oss:20b} & plain & 0.93 & 0.96 & 0.83 & 2 \\
\texttt{gpt-oss:20b} & fewshot & 1.00 & 1.00 & 1.00 & 0 \\
\bottomrule
\end{tabular}
\caption{Prompt-robustness ablation ($n=63$). The small-Gemma regression is a
\emph{capability} limit, not a format artefact: gemma3:4b emits clean parseable JSON
that verbalises the right distinction and still mislabels (manual inspection); switching to
a single-word format only shifts its threshold (FP $25\!\to\!11$); few-shot
demonstration partially rescues gemma3:12b but not gemma3:4b.}
\label{tab:prompt_robustness}
\end{table*}


\section{Realism: scoping to answer-bearing pairs}
\label{sec:realism}
A first realism run with HotpotQA-distractor produced an honest negative. Pairwise
conflict recall stayed low ($\le 0.45$) even after we fixed the refutation generator
(boolean-flip opposite-answer refutations). On this regime ($n=59$ pairs from the
corrected refutation generator), the cross-encoder hits F1 $0.18$, LLM-bare $0.62$,
and the question-conditioned judge $0.26$. Manual inspection clarified the cause. A
HotpotQA ``supporting'' paragraph is an \emph{evidence fragment about one entity}
(``Scott Derrickson is an American director''). An opposite-answer refutation is a
\emph{multi-hop relationship verdict} that often \emph{agrees on the shared fact}
(``born in Denver, Colorado'' is still American) and differs only on the combined
conclusion. S and R therefore \emph{do not contradict pairwise}. The judge's
\texttt{neutral} is correct, and the gold is ill-posed. Pairwise conflict detection is
only well-defined when \emph{both} passages assert an answer to the same question. We
accordingly scope the experiments to answer-bearing units.

\section{Long answer-bearing pairs}
\label{sec:answerbearing}
Restricting to 14 single-hop factoid questions, we LLM-write a 3 to 4 sentence passage
asserting each candidate answer and build conflict / agreement / topic-shift pairs over
the resulting long text (Table~\ref{tab:answer_bearing}). Conflict recall recovers
across detectors (cross-encoder $0.05 \to 0.50$ versus multi-hop evidence; LLM-bare
$0.45 \to 0.86$; answer-relevant $0.15 \to 0.79$). When the gold is well-posed and
both sides are answer-bearing, detectors detect. On long text, the LLM judges dominate
F1 ($0.89$ vs $0.67$). The cross-encoder regains perfect precision but length caps its
recall at $0.50$. Answer-relevance has perfect precision at a small recall cost. Its
big sense-trap precision win is muted here because these topic-shift distractors are
different-\emph{subject}, not same-name.

\begin{table}[t]
\centering
\small
\begin{tabular}{lccc}
\toprule
Detector & Precision (95\% CI) & Recall (95\% CI) & F1 \\
\midrule
cross\_encoder\_nli & 1.00 [1.00,1.00] & 0.50 [0.25,0.78] & 0.67 \\
llm\_bare & 0.92 [0.75,1.00] & 0.86 [0.65,1.00] & 0.89 \\
llm\_answer\_relevant & 1.00 [1.00,1.00] & 0.79 [0.55,1.00] & 0.88 \\
\bottomrule
\end{tabular}
\caption{Long answer-bearing passage pairs ($n=42$:
14 conflict / 14 agreement /
14 topic-shift, all 3--4 sentence LLM-written passages asserting a
specific answer). LLM judges dominate F1; the cross-encoder regains perfect precision but
length caps recall.}
\label{tab:answer_bearing}
\end{table}


\section{External validation on real FEVER}
\label{sec:fever}
We sanity-check the previous result on real human-labelled data.
\texttt{copenlu/fever\_gold\_evidence} bundles claim + label (SUPPORTS / REFUTES) +
gold evidence sentences. We evaluate the three detectors on 100 VERIFIABLE pairs
balanced 50/50 (Table~\ref{tab:fever}). The LLM judges keep \textbf{perfect precision}
($1.00$~[$1.00,1.00$]) on real data, so the precision finding survives external
validation. The cross-encoder, on its native training distribution (short
claim/evidence NLI), is \emph{competitive on F1} ($0.85$ vs $0.84/0.81$) via higher
recall ($0.80$ vs $0.72/0.68$). Answer-relevance neither helps nor hurts here (no
sense-trap failures to fix on FEVER pairs), consistent with the answer-bearing result.

\begin{table}[t]
\centering
\small
\begin{tabular}{lccc}
\toprule
Detector & Precision (95\% CI) & Recall (95\% CI) & F1 \\
\midrule
cross\_encoder\_nli & 0.91 [0.82,0.98] & 0.80 [0.68,0.91] & 0.85 \\
llm\_bare & 1.00 [1.00,1.00] & 0.72 [0.59,0.84] & 0.84 \\
llm\_answer\_relevant & 1.00 [1.00,1.00] & 0.68 [0.53,0.81] & 0.81 \\
\bottomrule
\end{tabular}
\caption{External validation on real human-labelled FEVER pairs
($n=100$; \texttt{copenlu/fever\_gold\_evidence}). On short claim--evidence NLI, the
cross-encoder's native distribution, it is competitive; LLM judges retain perfect precision
(no false positives on real data) but are more conservative on recall.}
\label{tab:fever}
\end{table}


\section{Vector-retrieval realism}
\label{sec:retrieval}
The experiments above evaluate the three detectors on curated or pre-paired data. This
section closes the deployment loop. We embed all 100 FEVER gold evidence sentences with
\texttt{all-MiniLM-L6-v2} and use exact cosine retrieval as a deterministic small-corpus
proxy for the vector-index stage, sidestepping ANN variance in this audit. We retrieve
top-$k=5$ passages for each of the 100 claims. Each
retrieved (claim, passage) pair is then scored by all three detectors. This is the
pipeline a vector-DB practitioner actually runs: dense top-$k$ followed by pairwise
conflict-checking on the returned passages.

Retrieval works as expected. Gold evidence is in the top-5 for $0.87$ of claims
(recall@5) and at rank 1 for $0.76$ (recall@1). The detector behaviour on the
\emph{retrieved} pairs is the interesting part (Table~\ref{tab:retrieval_realism}).
The cross-encoder flags conflict on $0.43$ [$0.39, 0.48$] of retrieved pairs. The two
LLM judges flag $0.09$ and $0.08$. Most retrieved passages for a given claim are
topically related but not contradicting (the same Wikipedia article cluster), so a
high flag rate here is a potential over-flagging surface in deployment, not a true
detection win.
On the subset where the retrieved passage is the FEVER gold evidence (where the
SUPPORTS/REFUTES label gives a per-pair ground truth), all three detectors are within
a few points of each other on precision-on-gold ($0.85$, $0.93$, $0.89$), so the
cross-encoder's extra flags are coming from the non-gold retrieved pairs, not from a
better gold signal. Latency separates the detectors by roughly $30\times$
($0.16$\,s for the cross-encoder, $4$ to $5$\,s for the LLM judges). The cost ladder
and the over-flagging ladder both push the same way on long real retrieved pairs as
they did on the synthetic answer-bearing pairs (\S\ref{sec:answerbearing}) and the
FEVER claim/evidence pairs (\S\ref{sec:fever}).

\begin{table}[t]
\centering
\small
\begin{tabular}{lccc}
\toprule
Detector & Flag rate (95\% CI) & Prec@gold & Latency (s) \\
\midrule
cross-enc NLI & 0.43 [0.39,0.48] & 0.85 & 0.16 \\
LLM bare & 0.09 [0.06,0.11] & 0.93 & 4.05 \\
LLM answer-rel & 0.08 [0.06,0.11] & 0.89 & 5.21 \\
\bottomrule
\end{tabular}
\caption{Vector-retrieval realism. Top-$k$ dense retrieval ($k=5$, exact cosine over
an embedded corpus of $n=100$ FEVER evidence sentences, all-MiniLM-L6-v2
embeddings) for $n=100$ FEVER claims; gold-evidence recall@$k$ is
0.87 (recall@1 $0.76$). Each retrieved (claim, passage) pair is
scored by the three detectors. The cross-encoder's higher flag rate at comparable
precision-on-gold reflects the same precision/recall trade seen on synthetic and FEVER
pairs (\S\ref{sec:answerbearing}, \S\ref{sec:fever}). Per-detector mean latency
shows the cost ladder a vector-DB practitioner faces in deployment.}
\label{tab:retrieval_realism}
\end{table}


\section{Synthesis: an input-regime ladder}
\label{sec:synthesis}
The detector choice is not a single ladder but an \emph{input-regime} choice
(Table~\ref{tab:synthesis}). The LLM-judge advantage is conditional on
(a)~the presence of same-name / sense-trap failures (where answer-relevance dominates)
and (b)~passage length (where the cross-encoder breaks down). On the cross-encoder's
short native NLI distribution, it is competitive. The practical recommendation: use a
question-conditioned LLM judge on \emph{answer-bearing} units (single-hop / claim-level).
The cross-encoder is appropriate only for short native-distribution pairs where high
recall matters more than the sense-trap precision risk.

\begin{table}[t]
\centering
\footnotesize
\begin{tabular}{@{}p{0.28\linewidth}p{0.32\linewidth}p{0.26\linewidth}@{}}
\toprule
Input regime & Best detector & Why \\
\midrule
Short curated, sense traps present (\S\ref{sec:answerrel}) &
  \textsc{llm\_answer\_relevant} (P $1.00$) & only it is question-conditioned \\
Short real claim--evidence (FEVER, \S\ref{sec:fever}) &
  cross-enc $\approx$ \textsc{llm\_bare} on F1 & cross-encoder's home turf \\
Real top-$k$ retrieved pairs (\S\ref{sec:retrieval}) &
  LLM judges (flag $0.09$ vs $0.43$) & cross-enc over-fires on non-gold \\
Long answer-bearing passages (\S\ref{sec:answerbearing}) &
  LLM judges (F1 $0.89$) & cross-encoder length-capped \\
Long multi-hop fragments (\S\ref{sec:realism}) &
  ill-posed; scope to answer-bearing & pairwise conflict undefined \\
\bottomrule
\end{tabular}
\caption{Detector choice depends on input regime. The LLM advantage is conditional on
sense-trap presence and passage length; the cross-encoder remains useful for short
native-distribution NLI.}
\label{tab:synthesis}
\end{table}

\paragraph{Implications for vector-DB practitioners}
A practitioner building a conflict-aware RAG pipeline on top of a vector store typically
faces three concrete choices: which detector to deploy, on what unit of retrieved
content, and how to interpret the resulting signal under model heterogeneity. The
input-regime ladder (Table~\ref{tab:synthesis}) gives a direct recommendation for each.
For pipelines that retrieve and conflict-check \emph{short claim/evidence pairs} (e.g.,
fact-checking RAG and similar pairwise filters), the cross-encoder is the right
default. It matches the LLM judge on F1 with much lower compute, and the sense-trap
precision risk is small at sentence scale. For pipelines that retrieve
\emph{paragraph-length} passages directly into the conflict-detection stage, the
cross-encoder fails (F1 $0.18$ at our long-passage scale) and an LLM judge is required.
Question-conditioning the LLM judge is then the precision-maximising default, unless
the deployed LLM is small enough to fall into the capability regime
(Table~\ref{tab:prompt_robustness}). For pipelines that retrieve \emph{multi-hop
evidence fragments} (HotpotQA-style or graph-RAG-style sub-evidence), pairwise conflict
detection is ill-posed on the retrieved units themselves. Conflict-aware logic should
instead be deferred to a downstream synthesis step that operates on answer-bearing
representations (\S\ref{sec:realism}).

\section{Limitations}
\label{sec:limitations}
(1)~Single-author curated probes; no independent gold validation or inter-annotator
agreement yet. (2)~Curated set is $n{=}63$. The orderings are statistically clean but
absolute gaps need a larger probe set (and a held-out probe author) to tighten.
(3)~The answer-bearing passages are LLM-generated. The FEVER validation is the
real-corpus counterfactual but covers only short claim/evidence pairs, not long
answer-bearing text on real corpora. (4)~The cross-encoder is fixed at
DeBERTa-base-MNLI as the canonical representative of the pairwise NLI baseline used
across the conflict-aware RAG
literature~\citep{conflictrag2025,madamrag2025,contradetect2025}. Larger NLI models may
shift the input-regime ladder. (5)~The 10-model sweep uses one prompt per condition;
the prompt-robustness study (\S\ref{sec:promptrob}) covers only three variants on
three models. (6)~We measure detector behaviour, not its downstream effect on RAG
generation. A follow-up should propagate the sense-trap precision win into end-task
faithfulness and abstention numbers. (7)~Query-conditioned NLI was concurrently
formalised as a benchmark by~\citet{canby2025qnli}. Our contribution is the deployment
audit, not the original technique; the audit is complementary to and extends their
benchmark.

\section{Conclusion}
\label{sec:conclusion}
We audit conflict detection on top of vector retrieval for RAG. Pairwise NLI is a
brittle signal. A same-name sense trap systematically false-positives, and the failure
is mechanistic (a missing question slot), not a model-capacity issue. The
question-conditioned formulation of~\citet{canby2025qnli} closes the gap when deployed
as an LLM-judge detector. We show the closure replicates across 8 of 10 open-weight
families with non-overlapping CIs. We also show that the choice between the pairwise
NLI cross-encoder commonly used in conflict-aware
RAG~\citep{conflictrag2025,madamrag2025,contradetect2025} and an LLM judge is
\emph{not} a global recommendation but an input-regime one. The cross-encoder wins on
short claim/evidence pairs and is unusable on long passages. LLM judges are robust
across lengths but conservative on recall. Small models cannot follow the
question-conditioning prompt and regress. Multi-hop evidence fragments are simply
ill-posed for pairwise judging. A small real top-$k$ retrieval check (over an embedded
FEVER evidence corpus) reproduces the same cost and over-flagging ladder. We package
the findings as a five-row deployment ladder for vector-DB-backed RAG practitioners,
with bootstrap confidence intervals throughout. Code and result artefacts are released.

\bibliographystyle{ACM-Reference-Format}
\bibliography{../references}
\end{document}
